% Generated draft: requires mathematical and source review.
\documentclass[11pt]{article}
\usepackage[margin=25mm]{geometry}
\usepackage{fontspec}
\setmainfont{Latin Modern Roman}
\usepackage{amsmath,amssymb,mathtools,unicode-math}
\setmathfont{Latin Modern Math}
\usepackage{enumitem,needspace,xurl,hyperref}
\setlength{\parindent}{0pt}
\setlength{\parskip}{5pt}
\setlength{\emergencystretch}{3em}
\newcommand{\N}{\mathbb{N}}
\newcommand{\Z}{\mathbb{Z}}
\newcommand{\Q}{\mathbb{Q}}
\newcommand{\R}{\mathbb{R}}
\newcommand{\C}{\mathbb{C}}
\newcommand{\F}{\mathbb{F}}
\newcommand{\A}{\mathbb{A}}
\newcommand{\PP}{\mathbb{P}}
\newcommand{\E}{\mathbb{E}}
\newcommand{\eps}{\varepsilon}
\newcommand{\dd}{\,\mathrm{d}}
\begin{document}
\section*{UnsolvedMath ID: 11100031}
\textbf{Source problem code:} catalogue entry
\subsection{Definitions and mathematical statement}
For a finite group $G$ and field $k$, let $\operatorname{StMod}(kG)$ be the stable category of all $kG$-modules, with morphisms taken modulo those factoring through projective modules. A localizing subcategory is a full triangulated subcategory closed under arbitrary set-indexed coproducts. It is a tensor ideal if $M\in\mathcal L$ implies $M\otimes_k N\in\mathcal L$ for every module $N$, with diagonal $G$ action.

Hovey's Applications item 7 asks to classify localizing subcategories and predicts a bijection with arbitrary subsets of $\operatorname{Proj}H^*(G;k)$, the homogeneous prime spectrum relevant to cohomological support. Schematically, the proposed support correspondence is
\[
 \mathcal L\longmapsto\bigcup_{M\in\mathcal L}\operatorname{supp}(M),
 \qquad
 S\longmapsto\{M:\operatorname{supp}(M)\subset S\}.
\]
The subsets are arbitrary, not required to be specialization closed.

A crucial qualification is absent from the excerpt and the registry's solution summary: for general finite groups the standard support classification is of tensor-ideal localizing subcategories. This does not automatically classify all localizing subcategories. For a finite $p$-group in characteristic $p$, the trivial module generates the stable category and the tensor-ideal condition is automatic; for general finite groups block and generation issues matter. The preceding source assumptions on $G$ must therefore be recovered before identifying the original unqualified request with that theorem.

The registry marks the problem solved by Benson--Iyengar--Krause stratification and claims all finite groups over arbitrary fields. That report is retained with the tensor-ideal qualification unresolved. If $kG$ is semisimple the stable category is zero, giving the corresponding trivial case. The all-module stable category must not be confused with its finite-dimensional compact subcategory and the separate thick-subcategory question.
\subsection{Short English statement}
Classify localizing stable-module subcategories by arbitrary cohomological-support subsets. The supplied solution claim needs a tensor-ideal qualification for general finite groups, or the original finite-$p$-group hypothesis.
\subsection{Sources}
\begin{itemize}
\item Local registry: \url{data/unsolvedmath/problems.json}, permanent UnsolvedMath ID $11100031$, exact source code \texttt{catalogue entry}, attributed to Mark Hovey.
\item \emph{Hovey's Algebraic Topology Problem List}, Applications item 7: \url{https://www-users.cse.umn.edu/~tlawson/hovey/}. Supplied access date: 2026-07-29; extraction: html-ordered-list; 65 top-level items in nine category pages.
\item Supplied solution attributions: Benson--Iyengar--Krause, ``Colocalizing subcategories and cohomological support'' and ``Localizing subcategories of the stable module category'' (given as 2018), with their Hopkins--Neeman-style stratification framework. Exact bibliographic matching, dates and hypotheses require review.
\item The registry compares the thick-subcategory item with source code \texttt{AMR-110-0030}; that separate record's definitions are not supplied here.
\end{itemize}
\end{document}
